\section{Standard and nonstandard models of arithmetic}
\label{sec:arithmetic-coding-02}
\label{sec:nonstandard-arithmetic}

\begin{definition}[Standard and nonstandard models]
\label{def:nonstandard-model}
Let \(T\) be an \(\Lang_A\)-theory satisfied by \(\mathcal N\).  A model
\(\mathcal M\) of \(T\) is \emph{standard} iff
\(\mathcal M\cong\mathcal N\) as an \(\Lang_A\)-structure, and it is
\emph{nonstandard} otherwise.  These terms classify models of arithmetic by
isomorphism type; they do not classify sentences or proof predicates.  The
distinction is stronger than agreement on all sentences, as the theorem and
exercise below show, and it is needed to state what compactness produces.
\end{definition}

\begin{theorem}[Compactness produces nonstandard models]
\label{thm:nonstandard-model}
Assume the Axiom of Choice in the ambient set theory.  Every \(\Lang_A\)-theory
\(T\) satisfied by \(\mathcal N\) has a countable
nonstandard model.
\textit{Status: complete proof.}
\end{theorem}

\begin{proof}
Add one new constant \(c\) and consider
\[
  T\cup\{c\ne\bar n:n\in\N\}.
\]
Every finite subset has an expansion of \(\mathcal N\) as a model: interpret
\(c\) by a natural number different from the finitely many numerals mentioned.
By first-order compactness, the whole set has a model \(\mathcal M'\).  In its
\(\Lang_A\)-reduct \(\mathcal M\), the element \(c^{\mathcal M'}\) is unequal
to \(\bar n^{\mathcal M}\) for every standard \(n\).  If
\(\mathcal M\cong\mathcal N\), every element of \(M\) would be the value of
some numeral, a contradiction.  Thus the reduct is a nonstandard model of
\(T\).  Applying \cref{thm:downward-ls} in the expanded countable signature,
with \(c^{\mathcal M'}\) among the designated elements, makes the model
countable while preserving all the displayed sentences.
\end{proof}

See \cite{Hodges1993,Marker2002} for further compactness constructions of
nonstandard models.

\begin{example}[Nonstandard Peano arithmetic]
Since \(\mathcal N\entails\PA\), \cref{thm:nonstandard-model} gives a
countable model of \(\PA\) containing an element unequal to every numeral.
The theorem is an existence result: it does not identify that element with
the value of any numeral, and it does not make the resulting model
isomorphic to
\(\mathcal N\).
\end{example}

\begin{exercise}
\label{ex:arithmetic-coding-01}
\label{ex:nonstandard-true-arithmetic}
Apply \cref{thm:nonstandard-model} to \(T=\Th(\mathcal N)\).  Prove that the
resulting countable model agrees with \(\mathcal N\) on every
\(\Lang_A\)-sentence but is not isomorphic to \(\mathcal N\).  Explain why
elementary equivalence therefore does not imply isomorphism.
\end{exercise}

\begin{warning}[Standard truth versus formal proof]
The notation \(\mathcal N\entails\sigma\) says that \(\sigma\) is true in the
standard structure.  The notation \(T\derives\sigma\) says that a finite
derivation exists from the axioms of \(T\) in a fixed effective calculus.
Passage from \(T\derives\sigma\) to \(\mathcal N\entails\sigma\)
requires soundness of the fixed calculus together with the hypothesis
\(\mathcal N\entails T\).  Passage from \(\mathcal N\entails\sigma\) to
\(T\derives\sigma\) fails in general even though the fixed calculus is
complete---this failure is the content of the first incompleteness theorem
below---and holds only for restricted sentence classes, such as
\(\Sigma^0_1\) sentences over theories extending \(\mathsf Q\).
\end{warning}

